\section{Scope, Corpus, and Qualifications}
\label{app:scope}

\subsection{Purpose, Intended Reader, and Organizing Thesis}

This reference work presents Catholic doctrine on the angels from within the Church's living tradition, tracing how Sacred Scripture, the Church Fathers, the \emph{Celestial Hierarchy} of Pseudo-Dionysius, and St.~Thomas Aquinas fit together into one coherent account of spiritual creation. It is written for serious Catholic readers, seminarians, catechists, and theological students willing to engage philosophical and patristic precision supported by primary citations. Its organizing thesis is that angelic excellence is wholly received and ordered to Christ: pure immateriality is not divinity, natural intellectual brilliance is not the beatific vision, hierarchical mediation does not rival God's sovereignty, and holy angels and redeemed human beings form one City of God under Jesus Christ, their Head.

\subsection{Corpus and Completeness}

The systematic spine consists of all twenty-four questions and all 118~articles of St.~Thomas Aquinas's treatise on the angels in the \emph{Summa theologiae}: Prima Pars questions~50--64 (72~articles) and questions~106--114 (46~articles), supplemented by Tertia Pars q.~8, a.~4 (Christ's headship over the angels) and q.~30, aa.~1--4 (the Annunciation), together with directly inspected parallels in \emph{De malo} q.~16 (all 12~articles), \emph{De veritate} qq.~8--9, \emph{De ente et essentia} chs.~4--5, and the Pauline commentaries on Ephesians and Colossians (read in the 1857 Dessain printing as a finding aid). Completeness means that every single one of those 118~articles is expounded in the body and listed exactly once in Appendix~\ref{app:article-census}.

All fifteen chapters of Pseudo-Dionysius's \emph{Celestial Hierarchy} are treated in order, with chapter titles from John Parker's 1899 translation, Migne \emph{Patrologia Graeca}~3 column bounds verified directly in the Greek-Latin edition, focused quotations with section numbers, and exact \emph{Summa} reception citations. Pre-Christian writings (Septuagint distinctive readings, \emph{1~Enoch}, \emph{Jubilees}, \emph{Testaments of the Twelve Patriarchs}, Philo, Hesiod, Plato, Plutarch, Apuleius, Plotinus, Porphyry, and Iamblichus), patristic witnesses, magisterial acts, later Catholic schools, Byzantine theologians, Roman and Byzantine liturgical calendars, and saints' writings are treated selectively at identified primary loci.

\subsection{Grading Rule}

Authority grades attach to individual theological propositions rather than to authors as a whole, using seven reader-facing categories: \textbf{Defined} (conciliar definition or solemn creed), \textbf{Church teaching} (ordinary magisterium, catechisms, papal acts, or universal liturgical confession), \textbf{Common teaching} (held commonly by Fathers and theologians), \textbf{Thomist position} (Aquinas's distinctive systematic determination within Catholic orthodoxy, naming contrary Catholic schools where documented), \textbf{Disputed} (an actual open disagreement among identified Fathers or schools without a prevailing common resolution), \textbf{Pious tradition} (devotional practice and saintly accounts received in Catholic spirituality), and \textbf{Apocryphal} (extra-canonical Jewish writings outside the biblical canon). Grade assignments represent source-grounded editorial synthesis; where witnesses leave a point unsettled, the lower grade is assigned.

\subsection{Terminology, Translation, Rights, and Witness Ceilings}

English Scripture quotations reproduce the public-domain Douay-Rheims (Challoner) text using Vulgate book and verse numbering, with Hebrew or modern Greek numbering added in brackets where they differ (for example, Ps~90:11 [Heb~91:11] and 1~Thess~4:15 [4:16]). Septuagint quotations follow Sir Lancelot Brenton's public-domain 1851 English translation, with Greek Psalms checked in H.~B.~Swete's text. Greek terms are transliterated in italic (\gk{angelos}, \gk{Dynameis}, \gk{Exousiai}) for consistent rendering across print and web surfaces; Appendix~\ref{app:terminology-witnesses} records translation collisions, especially between Latin scholastic choir names and Parker's 1899 rendering of the middle triad.

Patristic and scholastic quotations are drawn from unrevised public-domain printings. Where an electronic edition carries modernized New Advent English revisions (Origen, \emph{Commentary on Matthew} XIII.26--29; Augustine, \emph{On the Trinity} III.4--9), the exposition paraphrases accurately rather than presenting modernized text as an unrevised nineteenth-century printing. Three electronic transcription defects in Project Gutenberg's English \emph{Summa} (I q.~50 a.~2 ``individual'' for ``intellectual''; misnumbered bracket headers at I q.~60 a.~3 and q.~62 a.~9) were verified and corrected against the 1889 Leonine Latin text and the 1920 English Dominican revision.

Several works discussed in the historical survey---including Aquinas's \emph{Summa contra gentiles}, \emph{De spiritualibus creaturis}, \emph{De substantiis separatis}, and \emph{In II Sent.}; full texts of \emph{1~Enoch}, \emph{Testaments of the Twelve Patriarchs}, Philo, Hesiod, Plutarch, Apuleius, Plotinus, Porphyry, and Iamblichus; direct standalone editions of Scotus's \emph{Ordinatio}, Su\'arez's \emph{De angelis}, Maximus the Confessor's \emph{Mystagogia}, Theodore the Studite, and St.~Gregory Palamas; the Council of Vienne's \emph{Ad nostrum qui} (1312); and restricted modern Vatican and ICEL texts---were not consulted here in full electronic editions. Those works are paraphrased only insofar as supported by verified patristic quotations (notably St.~Augustine's \emph{City of God} VIII--X and XV.23), verified passage records (Plato's \emph{Symposium} 202d--203a; \emph{Jubilees} 10:1--11), the 1885 Quaracchi \emph{Scholia}, Hefele's conciliar history, or established doctrinal summaries; no verbatim quotation is reconstructed from any work not directly inspected in full text.

Received vocal prayers printed for recitation (\emph{Angele Dei} from the 1910 \emph{Raccolta}, No.~296; Pope Leo~XIII's Prayer to St.~Michael from the 1945 Lasance \emph{New Roman Missal}, p.~797) reproduce identified public-domain witnesses verbatim without adaptation. Project-authored prose is licensed under CC~BY~4.0; incorporated Scripture, liturgical texts, prayers, and third-party translations retain their own independent status.

\subsection{Pastoral Boundary and Review Status}

Intended as a study aid for theological understanding and prayerful reading, the present treatise is not an official liturgical book, catechism, magisterial act, or manual of deliverance, and it offers no instructions for exorcism, methods for diagnosing particular cases of spiritual affliction, or judgments about named persons; rites, canonical discipline, and pastoral norms belong to the separate study of Catholic exorcism. The work has passed automated census-parity, source-binding, schema, metadata, house-voice, web-edition, and print-layout checks; no independent human theological review or ecclesiastical approval is claimed.
